%% Tracksuit sprites for the TikZ figures (2026-11-05). \input this inside a group: it
%% redefines \sgrobot / \sgrobotoff / \sgguardgun / \sgguardsign (cogsci2027/figpics.tex) to draw sprites from sprites/,
%% and the change ends with the group, so other figures keep the vector robots.
%% The placement follows LLM Squid Game Thesis/poster/figs/tracksuit/standalone_tracksuit.tex:
%% the standing sprite's top meets the robot's top (.37) and is \sgspriteH robot units tall;
%% other poses keep the same head size, feet on the same line, head centre on x.
%% Every line ends in % (and meta.tex is read with \endlinechar=-1): the figures are set inside
%% \resizebox, where a stray space would widen the box and shrink the drawing.
\makeatletter%
\edef\sg@savedendline{\the\endlinechar}%
\endlinechar=-1 %
\expandafter\def\csname sw@stand_blue\endcsname{2.6681}\expandafter\def\csname sh@stand_blue\endcsname{5.9565}\expandafter\def\csname sx@stand_blue\endcsname{-0.0103}
\expandafter\def\csname sw@stand_orange\endcsname{2.6681}\expandafter\def\csname sh@stand_orange\endcsname{5.9565}\expandafter\def\csname sx@stand_orange\endcsname{-0.0103}
\expandafter\def\csname sw@stand_green\endcsname{2.6681}\expandafter\def\csname sh@stand_green\endcsname{5.9565}\expandafter\def\csname sx@stand_green\endcsname{-0.0103}
\expandafter\def\csname sw@stand_purple\endcsname{2.6681}\expandafter\def\csname sh@stand_purple\endcsname{5.9565}\expandafter\def\csname sx@stand_purple\endcsname{-0.0103}
\expandafter\def\csname sw@give_blue\endcsname{3.2662}\expandafter\def\csname sh@give_blue\endcsname{6.0378}\expandafter\def\csname sx@give_blue\endcsname{-0.3024}
\expandafter\def\csname sw@ask_orange\endcsname{2.9726}\expandafter\def\csname sh@ask_orange\endcsname{6.0482}\expandafter\def\csname sx@ask_orange\endcsname{-0.2371}
\expandafter\def\csname sw@walk_purple\endcsname{3.4696}\expandafter\def\csname sh@walk_purple\endcsname{6.9726}\expandafter\def\csname sx@walk_purple\endcsname{0.0904}
\expandafter\def\csname sw@walk_blue\endcsname{3.4696}\expandafter\def\csname sh@walk_blue\endcsname{6.9726}\expandafter\def\csname sx@walk_blue\endcsname{0.0904}
\expandafter\def\csname sw@walkbag_blue\endcsname{3.3785}\expandafter\def\csname sh@walkbag_blue\endcsname{6.0979}\expandafter\def\csname sx@walkbag_blue\endcsname{-0.0620}
\expandafter\def\csname sw@receive_orange\endcsname{3.6100}\expandafter\def\csname sh@receive_orange\endcsname{6.7994}\expandafter\def\csname sx@receive_orange\endcsname{0.1294}
\expandafter\def\csname sw@reach_blue\endcsname{4.0502}\expandafter\def\csname sh@reach_blue\endcsname{5.5903}\expandafter\def\csname sx@reach_blue\endcsname{0.3625}
\expandafter\def\csname sw@climb_orange\endcsname{3.5646}\expandafter\def\csname sh@climb_orange\endcsname{6.5649}\expandafter\def\csname sx@climb_orange\endcsname{0.1736}
\expandafter\def\csname sw@sit_orange\endcsname{3.5924}\expandafter\def\csname sh@sit_orange\endcsname{5.5246}\expandafter\def\csname sx@sit_orange\endcsname{-0.0589}
\expandafter\def\csname sw@pull_blue\endcsname{4.4803}\expandafter\def\csname sh@pull_blue\endcsname{5.7356}\expandafter\def\csname sx@pull_blue\endcsname{-0.7355}
\expandafter\def\csname sw@pull_orange\endcsname{4.4803}\expandafter\def\csname sh@pull_orange\endcsname{5.7356}\expandafter\def\csname sx@pull_orange\endcsname{-0.7355}
\expandafter\def\csname sw@pull_green\endcsname{4.4803}\expandafter\def\csname sh@pull_green\endcsname{5.7356}\expandafter\def\csname sx@pull_green\endcsname{-0.7355}
\expandafter\def\csname sw@walkL_purple\endcsname{3.4696}\expandafter\def\csname sh@walkL_purple\endcsname{6.9726}\expandafter\def\csname sx@walkL_purple\endcsname{-0.0904}
\expandafter\def\csname sw@off\endcsname{2.6681}\expandafter\def\csname sh@off\endcsname{5.9565}\expandafter\def\csname sx@off\endcsname{-0.0103}
%
\endlinechar=\sg@savedendline\relax%
\def\sgspriteTop{.37}%
\def\sgspriteH{.86}%
\def\sgspritedir{cogsci2027/ko/figures/tracksuit_20261105/sprites/}%
\def\sgcol@sgblue{blue}\def\sgcol@sgorange{orange}\def\sgcol@sggreen{green}\def\sgcol@sgpurple{purple}%
\def\sgU{\fpeval{\sgspriteH/\csname sh@stand_blue\endcsname}}%
%% \sgpose[s]{(x,y)}{sprite name}{label}
\providecommand{\sgpose}{}%
\renewcommand{\sgpose}[4][1]{%
  \node[inner sep=0pt,anchor=south] at ($#2+(\fpeval{-#1*\sgU*\csname sx@#3\endcsname},\fpeval{#1*(\sgspriteTop-\sgspriteH)})$)
    {\includegraphics[height=\fpeval{#1*\sgU*\csname sh@#3\endcsname}cm]{\sgspritedir#3.png}};%
  \node[text=black,font=\sffamily\bfseries\fontsize{8}{8}\selectfont,
    inner sep=0pt,anchor=south] at ($#2+(0,#1*.39)$) {#4};}%
\renewcommand{\sgrobot}[4][1]{\sgpose[#1]{#2}{stand_\csname sgcol@#3\endcsname}{#4}}%
\renewcommand{\sgrobotoff}[3][1]{\sgpose[#1]{#2}{off}{#3}}%
%% Guards as sprites in the same style (2026-11-05, thesis poster commit 9959bbe:
%% guard/build_guards.py -> sprites/guard_{aim,sign}.png and guard_meta.tex, coordinates in units
%% of the sprite height). Same arguments as the figpics macros; the feet centre sits on (x,y).
\endlinechar=-1 %
\expandafter\def\csname gd@aim@w\endcsname{0.7168}
\expandafter\def\csname gd@aim@fx\endcsname{0.4510}
\expandafter\def\csname gd@aim@mx\endcsname{0.0008}
\expandafter\def\csname gd@aim@my\endcsname{0.3078}
\expandafter\def\csname gd@sign@w\endcsname{0.5851}
\expandafter\def\csname gd@sign@fx\endcsname{0.3067}
\expandafter\def\csname gd@sign@x0\endcsname{0.1715}
\expandafter\def\csname gd@sign@x1\endcsname{0.5763}
\expandafter\def\csname gd@sign@y0\endcsname{0.8305}
\expandafter\def\csname gd@sign@y1\endcsname{0.9919}
\expandafter\def\csname gd@big@w\endcsname{0.6106}
\expandafter\def\csname gd@big@fx\endcsname{0.2844}
\expandafter\def\csname gd@big@hx0\endcsname{0.0660}
\expandafter\def\csname gd@big@hx1\endcsname{0.5007}
\expandafter\def\csname gd@big@hy1\endcsname{0.9879}
\expandafter\def\csname gd@big@mx\endcsname{0.2825}
\expandafter\def\csname gd@big@my\endcsname{0.7325}
%
\endlinechar=\sg@savedendline\relax%
\def\gd#1#2{\csname gd@#1@#2\endcsname}%
\def\sgguardH{1.15}% standing guard height (cm at scale 1)
\definecolor{sgboard}{HTML}{FFFFFF}%
\definecolor{sgboardline}{HTML}{1A1A1E}%
%% \sgguardgun[s]{(x,y)}{angle}: leaves (sgmuzzle); the sprite's gun angle is fixed, so angle is unused.
\renewcommand{\sgguardgun}[3][1]{%
  \pgfmathsetmacro\gdH{#1*\sgguardH}%
  \node[inner sep=0pt,anchor=south west] (gdimg) at ($#2+(\fpeval{-\gdH*\gd{aim}{fx}},0)$)
    {\includegraphics[height=\gdH cm]{\sgspritedir guard_aim.png}};%
  \coordinate (sgmuzzle) at ($(gdimg.south west)+(\fpeval{\gdH*\gd{aim}{mx}},\fpeval{\gdH*\gd{aim}{my}})$);}%
%% \sgguardsign[s]{(x,y)}{h}{text}: the board top meets the old text box top (origin + s(h+.42));
%% a vector board of the same shape covers the drawn white board, so long text only widens it.
\renewcommand{\sgguardsign}[4][1]{%
  \pgfmathsetmacro\gdH{#1*(#3+.42)/\gd{sign}{y1}}%
  \node[inner sep=0pt,anchor=south west] (gdimg) at ($#2+(\fpeval{-\gdH*\gd{sign}{fx}},0)$)
    {\includegraphics[height=\gdH cm]{\sgspritedir guard_sign.png}};%
  \node[fill=sgboard,draw=sgboardline,line width=.9pt,inner sep=2pt,
    minimum width=\fpeval{\gdH*(\gd{sign}{x1}-\gd{sign}{x0})+.04}cm,
    minimum height=\fpeval{\gdH*(\gd{sign}{y1}-\gd{sign}{y0})+.04}cm,
    font=\sffamily\bfseries\fontsize{7.4}{8}\selectfont,text=sgastra]
    at ($(gdimg.south west)+(\fpeval{\gdH*(\gd{sign}{x0}+\gd{sign}{x1})/2},\fpeval{\gdH*(\gd{sign}{y0}+\gd{sign}{y1})/2})$) {#4};}%
\makeatother%
